\DocumentMetadata{
  pdfversion=2.0,
  pdfstandard=ua-2,
  lang=en-US,
  tagging=on,
  tagging-setup={math/setup=mathml-SE,tabsorder=structure}
}
\documentclass[11pt,letterpaper]{article}
\usepackage[margin=0.68in,includefoot]{geometry}
\usepackage{newcomputermodern}
\usepackage{amsmath,amscd,mathtools}
\usepackage{tikz,adjustbox}
\providecommand{\square}{\mdlgwhtsquare}
\usepackage[hidelinks]{hyperref}
\hypersetup{
  pdftitle={Logic PhD Exam, May 2022},
  pdfauthor={University of Florida Department of Mathematics},
  pdfsubject={Graduate mathematics examination},
  pdfdisplaydoctitle=true,
  bookmarksopen=true
}
\setcounter{secnumdepth}{0}
\setlength{\parindent}{0pt}
\setlength{\parskip}{0.28em}
\setlength{\emergencystretch}{3em}
\setlength{\leftmargini}{2.15em}
\setlength{\leftmarginii}{2.65em}
\setlength{\leftmarginiii}{3em}
\pagestyle{empty}
\raggedbottom
\allowdisplaybreaks
\NewTaggingSocketPlug{block/list/label}{examproblem}{%
  \tagstructbegin{tag=Lbl}\tagstructend%
  \tagstructbegin{tag=\LBody}%
  \tagstructbegin{tag=\ProblemHeadingTag,title-o={\ProblemHeadingText}}%
  \tagmcbegin{tag=\ProblemHeadingTag,actualtext={\ProblemHeadingText}}%
  #2%
  \tagmcend%
  \tagstructend%
}
\newenvironment{examproblems}[2]{%
  \def\ProblemHeadingTag{#2}%
  \AssignTaggingSocketPlug{block/list/label}{examproblem}%
  \begin{enumerate}%
  \setcounter{enumi}{#1}%
  \renewcommand{\labelenumi}{\textbf{\theenumi.}}%
  \setlength{\topsep}{0.35em}%
  \setlength{\partopsep}{0pt}%
  \setlength{\itemsep}{0.48em}%
  \setlength{\parsep}{0.18em}%
}{\end{enumerate}}
\newenvironment{examsubparts}{%
  \AssignTaggingSocketPlug{block/list/label}{default}%
  \begin{enumerate}%
  \setlength{\topsep}{0.18em}%
  \setlength{\partopsep}{0pt}%
  \setlength{\itemsep}{0.16em}%
  \setlength{\parsep}{0.1em}%
}{\end{enumerate}}
\newenvironment{examinstructions}{%
  \par\begingroup\itshape
}{%
  \par\endgroup\vspace{0.18em}
}
\makeatletter
\renewcommand\subsection{\@startsection{subsection}{2}{0pt}%
  {-1.25ex plus -.25ex minus -.1ex}%
  {0.55ex plus .1ex}%
  {\normalfont\large\bfseries}}
\renewcommand{\@maketitle}{%
  \newpage\null\vskip -1em%
  {\footnotesize\bfseries
    \href{https://www.ufl.edu/}{University of Florida}, \href{https://math.ufl.edu/}{Department of Mathematics}\par}%
  \vskip -0.5em%
  \tagstructbegin{tag=H1,title={Logic PhD Exam, May 2022}}%
  \tagpdfparaOff%
  \tagmcbegin{tag=H1}%
    {\LARGE\bfseries\href{https://gma.math.ufl.edu/exams/logic-phd/}{Logic PhD Exam}, May 2022\par}%
  \tagmcend%
  \tagpdfparaOn%
  \tagstructend%
  \vskip 0.25em%
  \tagmcbegin{artifact}%
    \hrule height 1pt%
  \tagmcend%
  \vskip 1em%
}
\makeatother
\title{Logic PhD Exam, May 2022}
\author{}
\date{}
\begin{document}
\maketitle
\thispagestyle{empty}
\begin{examinstructions}
Solve 5 problems of the following; at least one from each section.
\end{examinstructions}
\subsection{A. Set Theory.}
\begin{examproblems}{0}{H3}
\def\ProblemHeadingText{Problem 1.}
\item
Sketch any proof of consistency of the Continuum Hypothesis with the axioms of ZFC.
\def\ProblemHeadingText{Problem 2.}
\item
What is a generic extension of a transitive model of ZFC? Provide a definition.
\def\ProblemHeadingText{Problem 3.}
\item
Show that every Polish space is a continuous image of the Baire space.
\end{examproblems}
\subsection{B. Computability.}
\begin{examproblems}{0}{H3}
\def\ProblemHeadingText{Problem 1.}
\item
Define many-one reducibility and many-one equivalence. Prove that there exist two non-computable, computably enumerable sets which are not many-one equivalent.
\def\ProblemHeadingText{Problem 2.}
\item
Prove Rice’s Theorem: If~\(P\) is a nonempty proper subset of the set of all partial computable functions, then the set of all indices of functions belonging to~\(P\) is not computable.
\def\ProblemHeadingText{Problem 3.}
\item
Define the Turing degrees and the Turing jump operator. State the Friedberg Jump Inversion Theorem and outline its proof.
\def\ProblemHeadingText{Problem 4.}
\item
State and prove Kleene’s Fixed Point Theorem (a.k.a. Recursion Theorem).
\end{examproblems}
\subsection{C. Model theory.}
\begin{examproblems}{0}{H3}
\def\ProblemHeadingText{Problem 1.}
\item
State both the upward and downward Löwenheim–Skolem Theorems. Prove either one of them.
\def\ProblemHeadingText{Problem 2.}
\item
Define what an ultraproduct of structures is and state Łoś’s Theorem. Use Łoś’s Theorem to give a proof of the Compactness Theorem.
\def\ProblemHeadingText{Problem 3.}
\item
Let \(\mathcal L\) be a language consisting of one unary function symbol~\(f\). Consider the class~\(K\) of all finite \(\mathcal L\)-structures in which~\(f\) is a bijection. Prove that~\(K\) is a Fraïssé class and describe its Fraïssé limit~\(M\). Prove or disprove: \(\operatorname{Th}(M)\) is \(\aleph_0\)-categorical.
\def\ProblemHeadingText{Problem 4.}
\item
Let \(\mathcal L\) be a language, let~\(n\) be a number, let~\(M\) be an \(\mathcal L\)-structure, and let~\(A\) be a subset of~\(M\). Give the definition of an~\(n\)-type of~\(M\) over~\(A\). Define the Stone topology on the space~\(S\) of complete~\(n\)-types of~\(M\) over~\(A\) and prove that~\(S\) equipped with this topology is a compact space.
\end{examproblems}
\end{document}
